% ===========================================================================
%  Beamer theme and slide helpers.
%
%  Loaded by the presentations (dokumente/trainer/) after packages.tex
%  and one of the command files -- never by the article documents, which use
%  style.tex instead.
%
%  It sets the colour theme, patches \section so the table of contents shows
%  page numbers, and provides the coloured full-slide macros the decks are
%  written with (\greenSlide, \blueSlide, ... -- see below).
% ===========================================================================

\newcommand{\newuncover}[2]{\uncover<#1>{#2}}
\newcommand{\term}[1]{\structure{#1}}
\newcommand{\halt}[0]{\pause}


% =========================================================
% Sectioning
% =========================================================


% =========================================================
% =========================================================
% include page numbers in toc:
\makeatletter
\long\def\beamer@section[#1]#2{%
  \beamer@savemode%
  \mode<all>%
  \ifbeamer@inlecture
    \refstepcounter{section}%
    \beamer@ifempty{#2}%
    {\long\def\secname{#1}\long\def\lastsection{#1}}%
    {\global\advance\beamer@tocsectionnumber by 1\relax%
      \long\def\secname{#2}%
      \long\def\lastsection{#1}%
      \addtocontents{toc}{\protect\beamer@sectionintoc{\the\c@section}{#2\hfill\the\c@page}{\the\c@page}{\the\c@part}%
        {\the\beamer@tocsectionnumber}}}%
    {\let\\=\relax\xdef\sectionlink{{Navigation\the\c@page}{\noexpand\secname}}}%
    \beamer@tempcount=\c@page\advance\beamer@tempcount by -1%
    \beamer@ifempty{#1}{}{%
      \addtocontents{nav}{\protect\headcommand{\protect\sectionentry{\the\c@section}{#1}{\the\c@page}{\secname}{\the\c@part}}}%
      \addtocontents{nav}{\protect\headcommand{\protect\beamer@sectionpages{\the\beamer@sectionstartpage}{\the\beamer@tempcount}}}%
      \addtocontents{nav}{\protect\headcommand{\protect\beamer@subsectionpages{\the\beamer@subsectionstartpage}{\the\beamer@tempcount}}}%
    }%
    \beamer@sectionstartpage=\c@page%
    \beamer@subsectionstartpage=\c@page%
    \def\insertsection{\expandafter\hyperlink\sectionlink}%
    \def\insertsubsection{}%
    \def\insertsubsubsection{}%
    \def\insertsectionhead{\hyperlink{Navigation\the\c@page}{#1}}%
    \def\insertsubsectionhead{}%
    \def\insertsubsubsectionhead{}%
    \def\lastsubsection{}%
    \Hy@writebookmark{\the\c@section}{\secname}{Outline\the\c@part.\the\c@section}{2}{toc}%
    \hyper@anchorstart{Outline\the\c@part.\the\c@section}\hyper@anchorend%
    \beamer@ifempty{#2}{\beamer@atbeginsections}{\beamer@atbeginsection}%
  \fi%
  \beamer@resumemode}%

\def\beamer@subsection[#1]#2{%
  \beamer@savemode%
  \mode<all>%
  \ifbeamer@inlecture%
    \refstepcounter{subsection}%
    \beamer@ifempty{#2}{\long\def\subsecname{#1}\long\def\lastsubsection{#1}}
    {%
      \long\def\subsecname{#2}%
      \long\def\lastsubsection{#1}%
      \addtocontents{toc}{\protect\beamer@subsectionintoc{\the\c@section}{\the\c@subsection}{#2\hfill\the\c@page}{\the\c@page}{\the\c@part}{\the\beamer@tocsectionnumber}}%
    }%
    \beamer@tempcount=\c@page\advance\beamer@tempcount by -1%
    \addtocontents{nav}{%
      \protect\headcommand{\protect\beamer@subsectionentry{\the\c@part}{\the\c@section}{\the\c@subsection}{\the\c@page}{\lastsubsection}}%
      \protect\headcommand{\protect\beamer@subsectionpages{\the\beamer@subsectionstartpage}{\the\beamer@tempcount}}%
    }%
    \beamer@subsectionstartpage=\c@page%
    \edef\subsectionlink{{Navigation\the\c@page}{\noexpand\subsecname}}%
    \def\insertsubsection{\expandafter\hyperlink\subsectionlink}%
    \def\insertsubsubsection{}%
    \def\insertsubsectionhead{\hyperlink{Navigation\the\c@page}{#1}}%
    \def\insertsubsubsectionhead{}%
    \Hy@writebookmark{\the\c@subsection}{#2}{Outline\the\c@part.\the\c@section.\the\c@subsection.\the\c@page}{3}{toc}%
    \hyper@anchorstart{Outline\the\c@part.\the\c@section.\the\c@subsection.\the\c@page}\hyper@anchorend%
    \beamer@ifempty{#2}{\beamer@atbeginsubsections}{\beamer@atbeginsubsection}%
  \fi%
  \beamer@resumemode}
\makeatother


% reduce spacing between toc items
\makeatletter
\patchcmd{\beamer@sectionintoc}{\vskip1.5em}{\vskip-0.1em}{}{}
\makeatother

\newcommand{\firsttoc}{
	\setbeamercolor{section in toc}{fg=yellow}
	\setbeamercolor{frametitle}{fg=white,bg=blue!20!white}
	\begin{frame}{\hyperlink{themenuebersicht}{Themenübersicht}}
		\label{themenuebersicht}
		\ \\
		\tableofcontents[
			sectionstyle=show/show,
			subsectionstyle=hide/hide/hide
		]
	\end{frame}
}

\newcommand{\partstoc}{
	\setbeamercolor{section in toc}{fg=white}
	\setbeamercolor{frametitle}{fg=white,bg=blue!20!white}
	\begin{frame}{\hyperlink{themenuebersicht}{Themenübersicht}}
		\ \\
		\tableofcontents[
		sectionstyle=show/shaded,
		subsectionstyle=show/shaded/hide
		]
	\end{frame}
}

\newcommand{\chapterstoc}{
	\setbeamercolor{section in toc}{fg=black}
	\setbeamercolor{frametitle}{fg=white,bg=black}
	\begin{frame}{\hyperlink{themenuebersicht}{Themenübersicht}}
		\ \\
		\tableofcontents[
			sectionstyle=show/hide,
			subsectionstyle=show/shaded/hide
		]
	\end{frame}
}

\newcommand{\newpart}[1]{
	\section{#1} 
	\partstoc
}

\newcommand{\newchapter}[1]{
	\subsection{#1}
	\chapterstoc
}

\newcommand{\newsection}[1]{}

% =========================================================
% Blocks
% =========================================================

\newcommand{\newslide}{\framebreak}

\newcommand{\createblock}[3]{
	\setbeamercolor{frametitle}{fg=#1,bg=#1!20}
	\setbeamercolor{structure}{bg=white,fg=#1}
	\setbeamercolor{block title}{fg=#1,bg=white}
	\setbeamercolor{itemize item}{fg=#1}
	\setbeamercolor{itemize subitem}{fg=#1}
	\setbeamercolor{enumerate item}{fg=#1}
	\setbeamercolor{enumerate subitem}{fg=#1}
	\setbeamertemplate{frametitle continuation}{}	% supress numbering I, II, III,...
	\begin{frame}[c,allowframebreaks]{#2}
		\small #3
	\end{frame}
}

% text
\newenvironment{mtext}[0]{\comment}{\endcomment}					% book
\NewEnviron{btext}[0]{\createblock{black}{}{\BODY}}					% beamer
\NewEnviron{atext}[0]{\createblock{black}{}{\BODY}}					% any

% blocks
\NewEnviron{dblock}[1]{\createblock{blue!75!black}{#1}{\BODY}}		% definition
\NewEnviron{tblock}[1]{\createblock{blue!75!black}{#1}{\BODY}}		% theorem
\NewEnviron{eblock}[1]{\createblock{green!50!black}{#1}{\BODY}}		% example
\NewEnviron{qblock}[1]{\createblock{gray!90!black}{#1}{\BODY}}		% quiz
\NewEnviron{ablock}[1]{\createblock{gray!90!black}{#1}{\BODY}}		% assignment
\NewEnviron{mblock}[1]{\createblock{gray!90!black}{#1}{\BODY}}		% model assignment
\NewEnviron{wblock}[1]{\createblock{red!95!black}{#1}{\BODY}}		% warning
\NewEnviron{iblock}[1]{\createblock{orange!90!yellow}{#1}{\BODY}}	% information
\NewEnviron{pblock}[1]{\createblock{blue!75!black}{#1}{\BODY}}		% proof
\NewEnviron{sblock}[1]{\createblock{black}{#1}{\BODY}}				% solution
\NewEnviron{triblock}[1]{\createblock{Fuchsia!90!black}{#1}{\BODY}}	% trivia

% =========================================================
% Footers
% =========================================================

\setbeamerfont{page number in head/foot}{size = \scriptsize}
\setbeamertemplate{footline}[frame number]
\setbeamertemplate{navigation symbols}{}


%----------------------------------------------------------------------------------------
%NEW:

\newcommand{\slide}[3]{
	\setbeamercolor{frametitle}{fg=black,#1}
	\begin{frame}{#2}
		#3
	\end{frame}
}
\newcommand{\greenSlide}[2]{
	\slide{bg=green!30!white}{#1}{#2}
}
\newcommand{\orangeSlide}[2]{
	\slide{bg=orange!50!white}{#1}{#2}
}
\newcommand{\graySlide}[2]{
	\slide{bg=gray!80!black}{#1}{#2}
}
\newcommand{\yellowSlide}[2]{
	\slide{bg=yellow}{#1}{#2}
}
\newcommand{\blueSlide}[2]{
	\slide{bg=blue!30!white}{#1}{#2}
}


\newcommand{\bsec}[2]{
	\section{#1}
	\partstoc
	\setbeamercolor{frametitle}{fg=black,bg=green!30!white}
	\begin{frame}{#1}
		#2
	\end{frame}
}
\newcommand{\Bsec}[2]{
	\section{#1}
	\partstoc
	\setbeamercolor{frametitle}{fg=black,bg=green!30!white}
	\begin{frame}{#1}
	#2% \sbreak
	\end{frame}
}

\newcommand{\bsubsec}[2]{
	\subsection{#1}
	\setbeamercolor{frametitle}{fg=orange!30!black,bg=orange!50!white}
	\begin{frame}{#1}
		#2
	\end{frame}
}

\newcommand{\btask}[1]{
	\setbeamercolor{frametitle}{fg=white,bg=gray!70!black}
	\begin{frame}{Aufgaben}
	#1
	\end{frame}
}
	
\newcommand{\bsol}[1]{
	\setbeamercolor{frametitle}{fg=black,bg=yellow}
	\begin{frame}{Lösungen}
	#1
	\end{frame}
}

\newcommand{\bquestion}[1]{
	\subsection{Verständnis Fragen}
	\setbeamercolor{frametitle}{fg=orange!30!black,bg=orange!50!white}
	\begin{frame}{Verständnis Fragen}
		Hier finden sich Zusatz Fragen aus dem letzten Oberkapitel, wo es möglichst nur um Verständnis geht.
		#1	
	\end{frame}
}

\newcommand{\tbd}{Hier gibt es noch nichts zu sehen!}

	%#1 }}
